\chapter{绪论}
\section{研究背景与意义}
\subsection{研究背景}
金融风险预测不仅涉及资产价格变化的方向，还涉及未来价格变动的幅度及其不确定性。相同的预期收益可以对应不同的风险水平，因而波动率预测是理解和管理风险的一项基础工作。高频价格所提供的日内变动信息，使研究者能够将一个交易日内部的价格路径转化为日度已实现测度，并在较低频率上开展预测研究。已实现测度的构造、条件波动的建模与预测分布之间的联系，构成了高频金融计量的重要研究基础\cite{andersen2003}。

已实现波动率的预测通常从资产自身历史出发。HAR 模型以日、周、月尺度的滞后波动率构造简约的预测关系，既便于估计，又为加入其他信息提供了清楚的参照\cite{corsi2009}。在多资产环境中，研究问题进一步扩展为：在某只股票的自身波动历史之外，其他股票当下的风险状态是否还包含有助于预测的信息？回答这一问题并不意味着应当舍弃自身历史，而是需要检验新增横截面信息能否在强基线之上提供增量。

横截面信息可以通过不同方式进入预测模型。市场或行业平均波动率提供总体状态，风险输入与输出指标概括一个节点在网络中的位置，逐边模型则保留不同邻居与目标股票之间的对应关系。三种表示保留的信息不同，复杂度也不同。本文关心的不是抽象意义上“使用网络是否更先进”，而是在风险已经可以被汇总为少数变量时，进一步保留每一条关系是否值得。

现有研究已经将图神经网络用于已实现波动率预测，也已将动态 Copula-CoVaR 生成的风险汇总指标引入 HAR 框架\cite{zhang2025,wang2026}。因此，本文不将“HAR 与 GNN 相结合”或“CoVaR 进入波动率预测”本身视为全新问题，而将研究重点放在关系信息的压缩与保留：当风险总量相近时，来自不同股票、具有不同时间变化的条件风险关系，是否应以不同方式参与目标股票的预测？

\subsection{理论意义}
本文的理论意义首先在于将风险度量与预测任务之间的联系具体化。CoVaR 是一个条件风险量，并不直接等于未来已实现方差。本文将其视为基于日收益条件分布得到的关系特征，检验其能否帮助预测另一个明确界定的对象，而不将“度量具有经济含义”作为“必然具有预测能力”的替代证据。

其次，本文区分关系表示的不同层次。二值图描述哪些股票相连，标量边描述关系强度，多维边进一步记录标准化强度及其变化。通过保持节点信息与骨干模型一致，可以将这种表示差异转化为可比较的模型限制，而不是一次性改变图、损失函数和模型容量之后，再笼统归因于风险网络。

最后，本文明确讨论风险语义与模型结构的边界。条件方向不等于因果方向；动态相关不等于上下尾不对称；未经标准化的双向 CoVaR 差异也可能主要体现目标股票波动尺度的不同。对这些概念作出区分，可以使研究结论与模型实际识别的对象保持一致。

\subsection{实践意义}
本文预期为风险预测提供一条可审计的实现路径。对于使用预测结果的研究人员，除了一个最终方差预测值，还需要知道输入数据的时段范围、风险边的估计截止日、图估计失败时的处理方式，以及模型在不同市场状态下是否稳定。将这些信息纳入实验记录，有助于区分模型误差、数据问题与使用时点问题。

实践价值需要建立在相应证据之上。预测误差的降低并不自动意味着交易收益提高，也不自动构成 VaR 或 ES 的有效性证明。本文首先关注预测任务本身；资产配置、波动率择时与风险管理应用只有在收益口径、信号可用时间和交易约束都可核验的情况下，才作为后续扩展分析。

\section{研究内容与研究问题}
\subsection{研究内容}
本文沿开题报告确定的研究路线\cite{proposal2026}，将工作组织为数据测量、风险关系估计、预测模型与样本外比较四个环节。第一，以公司分钟价格序列构造五分钟采样的日度已实现方差，明确集合竞价、午休、停牌、缺失和收益调整的处理。第二，用相同股票的日收益估计边际分布和动态依赖，从同一二元模型计算两个条件方向的风险量。第三，在自身 HAR 分支之外设置浅层逐边消息聚合，同时建立使用相同风险汇总信息的竞争模型。第四，按时间推进完成不同期限和训练损失的预测，结合关系破坏、尺度标准化和底层参数对照分析增量来源。

\subsection{研究问题}
本文围绕以下四个相互衔接的问题展开。第一，条件尾部风险关系在原始尺度与标准化尺度下呈现何种动态差异，其中多少可以由边际波动尺度或分布形状解释？第二，相比公共市场信息和风险输入、输出汇总，逐边关系是否提供额外的样本外预测价值？第三，在控制动态相关、边际风险与模型容量后，CoVaR 表征以及多维边条件化分别起什么作用？第四，上述价值是否随预测期限、市场状态、数据采样和估计窗口发生变化？

这里的“方向性”首先指条件设定中来源与目标的区分，不预设已识别真实市场中的单向传染。对问题的回答将依赖预先设定的比较和实际数据，而不依赖网络图的视觉复杂度。

\section{研究方法与技术路线}
\subsection{研究方法}
文献研究用于确定研究与已有工作的关系，尤其是图上的非线性溢出预测和 CoVaR 风险汇总预测之间的交叉。数学分析用于澄清条件分位数、对称性与尺度分解，使实现能够接受可复核的单元测试。实证比较则将不同信息集合与函数形式放入统一的时间外评价协议，分别检验其作用。

技术实现采用两条数据路径。分钟价格进入日度已实现方差及 HAR 特征，日收益进入动态条件风险边。两条路径在预测起点汇合，但未来真实收益与未来已实现方差不进入任何输入。风险图参数按预设频率重估，状态在可用观测到达后过滤，预测模型只用已经成熟的标签更新。

\subsection{论证顺序}
论文的论证从最简自身基线开始，依次增加公共信息、风险汇总、具体关系和多维边输入。这个顺序使每一步新增内容都有相应对照。若逐边模型只能优于未加入公共信息的 HAR，结论就限于横截面信息有效；只有进一步超过等信息风险汇总与动态相关等控制，才支持更强的关系级表征主张。

\section{研究设计的主要特点}
本文拟形成三个层面的研究贡献。其一，围绕“风险汇总是否压缩了对预测有用的关系”构建模型与实验，使研究动机直接对应可验证的问题。其二，将条件风险边的定义、标准化与归因控制纳入方法设计，明确风险尺度、相关参数和边表示各自的作用。其三，在公司高频数据场景中建立可追溯的数据与时间协议，使模型比较能够对应相同的样本、标签和信息集。

上述内容是本文的设计特点。对预测优越性、跨状态稳定性和应用价值的结论仍须由第五章的真实实证支持，不将方法组合或工作量本身视为已经证实的效果。

\section{论文结构}
全文分为六章。第一章提出研究背景、问题与总体路线。第二章梳理已实现波动率预测、动态条件风险及图预测文献。第三章给出研究对象、风险模型、边表示、预测表达式和必要数学性质。第四章说明公司数据合同、样本构造与实验协议。第五章按整体表现、贡献来源和稳健性组织实证结果。第六章根据证据总结结论并讨论实践启示与适用边界。
